\section{Related work}

The partially linear model originates in semiparametric regression, where a parametric coefficient is estimated alongside a nonparametric function of covariates \citep{Engle1986,Speckman1988}; \citet{Robinson1988} gave the root-\(n\) consistent residualized estimator that remains the reference construction, \citet{HardleMammen1993} developed the associated specification testing, and \citet{HardleLiangGao2000} give the monograph treatment. Partially linear models are a central workhorse for semiparametric estimation of low-dimensional structural coefficients in the presence of high-dimensional or otherwise flexible nuisance functions. The classical formulation and efficiency theory connect the coefficient \(\theta_0\), the partially linear coefficient, to residualized treatment variation after removing the treatment regression and outcome regression \citep{Hansen1982,Newey1990,Newey1994,NeweyMcFadden1994,BickelKlaassenRitovWellner1993,VanDerVaart1998,VanDerVaartWellner1996}. Modern double/debiased machine learning builds on the same orthogonalization principle by allowing flexible first-stage estimates while preserving first-order stability of the target moment \citep{BelloniChernozhukovHansen2014,ChernozhukovEtAl2018DML}. Within that tradition, the present paper focuses on a partially linear coefficient estimated from supplied treatment and outcome codes, and studies the additional identifying information carried by non-Gaussian treatment innovations.

A closely related line of work studies higher-order orthogonality and influence-function expansions for nuisance-robust semiparametric estimation. Higher-order constructions refine the first-order orthogonal-score idea by exploiting additional smoothness or distributional structure in nuisance errors \citep{RobinsRotnitzkyZhao1994,VanDerLaanRubin2006,RobinsLiTchetgenVanDerVaart2008,LiuMukherjeeNeweyRobins2017,NeweyRobins2018,LiuMukherjeeRobinsTchetgen2021,ChernozhukovEscancianoIchimuraNeweyRobins2022}. For the partially linear model, \citet{MackeySyrgkanisZadik2018} identify the role of non-Gaussian treatment residuals in higher-order orthogonal learning. \citet{JinMackeySyrgkanis2025} give the closest finite-order comparison point: their published order-\(r\) ACE guarantee, with ACE order \(r\), controls the \(1-\gamma\) error quantile, where \(\gamma\) is the probability level, by a constant multiple of
\[
\varepsilon_{1,n}^r\varepsilon_{2,n}
+
C_{\theta}\varepsilon_{1,n}^{r+1}
+
(\gamma n)^{-1/2},
\]
where \(\varepsilon_{1,n}\) and \(\varepsilon_{2,n}\) are the treatment- and outcome-code radii, \(C_{\theta}\) is the coefficient-range bound, and \(n\) is the sample size. The explicit prefactor \(r!\,16^r\delta^{-1}\), with \(\delta\) denoting the cumulant-separation level, makes the comparison sensitive to the finite order and to the amount of non-Gaussian separation.

The direct \(L^1(P_X)\) gate in the contour theorem plays a different operational role from the JMS pair of \(L^r(P_X)\) treatment and outcome radii. The contour gate controls the treatment-code perturbation of the residual transform and preserves the zero geometry; the ACE eligibility quantities in \cref{obj:def:jms-eligibility-quantities,obj:def:jms-eligibility} also use the outcome-code radius and the order-\(r\) finite-difference scale. On bounded covariate-regression ranges, an \(L^r(P_X)\) treatment radius at level \(\varepsilon_{1,n}\) gives the corresponding \(L^1(P_X)\) treatment control, so the common ACE class lies inside the contour class through \cref{obj:lem:jms-ace-class-relations}. The contour result therefore asks the treatment code to meet the displayed small-radius gate, while the published ACE guarantee asks the treatment and outcome radii to meet the JMS eligibility inequalities.

\begin{center}
\begin{tabular}{p{0.24\linewidth}p{0.33\linewidth}p{0.33\linewidth}}
\hline
Feature & Contour construction & JMS ACE comparison \\
\hline
Input accuracy & Current \(L^1(P_X)\) treatment-code radius and fixed primitive constants & Current \(L^r(P_X)\) treatment and outcome radii with JMS eligibility \\
Analytic resource & A localized zero of the treatment-innovation transform and a selected contour ratio & A finite order-\(r\) ACE expansion under cumulant separation \\
Risk criterion compared & Fixed-code MSE minimax bounds and a generalized-quantile upper bound for the contour statistic & Published generalized-quantile upper guarantee and class-level fixed-code MSE lower bound \\
Fixed-separation dependence & Constants depend on the fixed cumulant-separation level and primitive envelopes & The displayed imported bound carries the finite-order prefactor and \(\delta^{-1}\) dependence \\
\hline
\end{tabular}
\end{center}

This translation gives concrete nuisance-rate regimes for the comparison in \cref{obj:prop:jms-ace-alignment}. If \(\varepsilon_{1,n}=n^{-a}\) and \(\varepsilon_{2,n}=n^{-b}\), the ACE nuisance part dominates sampling when \(\min\{ra+b,(r+1)a\}<1/2\), subject to the JMS eligibility and contour small-radius gates; in that regime the ratio of the contour upper guarantee to the imported finite-order ACE upper guarantee tends to zero on the common clipped-code class. If \(\min\{ra+b,(r+1)a\}\ge1/2\), the ACE upper guarantee and the contour upper guarantee are both governed by the sampling term up to constants and eligibility factors. These comparisons concern the available upper bounds on the aligned class.

The contour approach developed here represents the same source of identifying variation through zeros of the treatment-innovation moment-generating function and converts those zeros into instruments and contour ratios. Transform methods have a long history in econometrics and statistics, including characteristic-function estimation, empirical transform methods, and moment-generating-function arguments \citep{FeuervergerMcDunnough1981,Singleton2001,ChackoViceira2003,CarrascoFlorensRenault2007,Carrasco2017}. The classical foundation for reading distributional structure off a transform is Marcinkiewicz's theorem, that a characteristic function of the form \(\exp(P)\) with \(P\) a polynomial forces \(\deg P\le2\) and hence a Gaussian law \citep{Marcinkiewicz1939}. In this paper's centered sub-Gaussian envelope, the fixed cumulant-separation condition yields the explicit zero-localization conclusion in \cref{obj:lem:zero-localization}: the treatment-innovation MGF has a zero inside the displayed disk. The zero-based step also connects to classical results on distributions, transforms, and completeness \citep{KaganLinnikRao1973,Mattner1992,DHaultfoeuille2011,DarollesFanFlorensRenault2011,HuShiu2022}. In the present setting, those tools are applied to residualized treatment variation formed with a supplied treatment code, so the analytic question becomes stability of transform zeros under \(L^1(P_X)\) treatment-code perturbations.

Non-Gaussianity is also a recurring identifying resource in independent component analysis, structural shock recovery, and related econometric models \citep{Comon1994,HyvarinenOja1997,HyvarinenKarhunenOja2001,LanneMeitzSaikkonen2017,LeeMesters2024,HoeschLeeMesters2024,ReizingerEtAl2025}. Relatedly, \citet{Andrews2017} exhibits families that are \(L^2\)-complete or boundedly complete, the property that makes a vanishing transform functional informative under the relevant class. The construction here uses non-Gaussianity through localized transform-zero geometry: it localizes a transform zero and integrates the observable ratio on a selected circle around the enclosed zero structure. The paper applies that resource in a partially linear estimation problem: cumulant separation supplies a transform zero, and the contour construction converts the resulting analytic geometry into a coefficient estimator. This places the contribution between higher-order semiparametric learning and non-Gaussian identification, with the comparison to ACE made on the common cumulant-separated class where both sets of assumptions are imposed.

The fixed-separation statistical statement is expressed in minimax terms. Under fixed cumulant separation and stable supplied treatment codes, the contour estimator attains the parametric \(n^{-1}\) mean-squared-error order over the non-Gaussian spectral class, while the comparison with \citet{JinMackeySyrgkanis2025} records the published ACE generalized-quantile upper guarantee on the aligned ACE class. The sequence comparison in \cref{obj:prop:jms-ace-alignment} is an upper-guarantee comparison on the common clipped-code class: when the displayed ACE nuisance terms dominate \(n^{-1/2}\), the ratio of the fixed-separation contour upper guarantee to the imported finite-order ACE upper guarantee tends to zero.

Finally, the finite-construction component draws on computable analysis and interval arithmetic. The relevant background treats real-number computation, constructive numerical representations, and certified interval enclosures \citep{Weihrauch2000,Ko1991,Moore1966}. Those tools provide the arithmetic substrate for the contour bank and represented-data statistic, while the statistical exposition keeps the identification and risk arguments in standard econometric form.
